% ============================================================================
% VI. KẾT LUẬN
% ============================================================================
\section{KẾT LUẬN}

\begin{itemize}
    \item \textbf{Kết quả thực hiện:} Báo cáo đã thiết kế và hiện thực hóa một kiến trúc RAG động định hướng \texttt{Coding Agent}, thay thế mô hình RAG tĩnh ở cả hai đầu ingestion và runtime:
    \begin{enumerate}
        \item Về \textbf{kiến trúc framework và bộ công cụ cho \texttt{Coding Agent}}, ở pha ingestion hệ thống thay thế \texttt{generic chunker} bằng \texttt{Parser interface} gồm 5 output invariants trên nền PostgreSQL 16 đa index (\texttt{pgvector} HNSW 512 chiều từ \texttt{Qwen3-Embedding-0.6B}, \texttt{ltree}, FTS \texttt{simple}, \texttt{pg\_trgm}) để \texttt{Coding Agent} tự viết \texttt{Parser} chuyên biệt cho từng miền tài liệu. Ở pha runtime, hệ thống thay thế lối \texttt{one-shot context stuffing} thụ động bằng bộ 6 MCP tools trên nền \texttt{Warm Streamable HTTP}, kết hợp \texttt{dual-channel retrieval engine} (Hybrid RRF + Cross-Encoder \texttt{mmarco-mMiniLMv2} hợp nhất với kênh regex/trigram theo \texttt{Venn partitioning}, gom nhóm theo \texttt{reading order} và đánh giá độ tin cậy qua \texttt{Sparse Match Gating}) với các công cụ duyệt cây \texttt{ltree}, đọc \texttt{raw line slice}, tự curate \texttt{Knowledge Graph} và kiểm chứng trích dẫn tất định trên PostgreSQL.
        \item Về \textbf{thử nghiệm lắp ghép \texttt{Parser} đa miền và điều tra \texttt{multi-hop}}, ba bộ \texttt{Parser} cho tài liệu kỹ thuật Next.js MDX, Node.js Markdown và Văn bản quy phạm pháp luật Bảo hiểm Việt Nam đã được lắp vào framework, indexing thành công 170 tài liệu gồm 171.183 dòng raw text thành 5.975 \texttt{chunk} bảo toàn 100\% code block, cấu trúc điều khoản và vị trí dòng 1:1. Kết quả đánh giá thực nghiệm trên 100 truy vấn thuộc 2 bộ test độc lập (\texttt{TS1-Lexical-API} và \texttt{TS2-Conceptual}) đạt tỷ lệ tìm thấy \textbf{Found = 97,0\%} (\textbf{100,0\%} trên toàn bộ 97/97 truy vấn có phạm vi lọc tài liệu hợp lệ và \textbf{100,0\%} trên toàn bộ miền Node.js), \textbf{Recall@5 = 92,5\%--93,0\%} (\textbf{95,0\%} trên \texttt{TS1} và \textbf{100,0\%} trên Node.js \texttt{TS1}), \textbf{Recall@10 = 95,0\%}, \textbf{MRR@10 = 78,9\%--79,0\%} (\textbf{83,4\%} trên \texttt{TS1}) và \textbf{NDCG@10 = 85,0\%} với độ trễ truy vấn trung vị $206,5\text{--}207,6\text{ ms}$, đồng thời giúp \texttt{Coding Agent} tự điều phối chuỗi công cụ giải quyết trọn vẹn các tình huống hỏi đáp pháp lý \texttt{multi-hop} có kiểm chứng trích dẫn xuyên suốt Luật 08/2022/QH15, Nghị định 46/2023/NĐ-CP và Thông tư 67/2023/TT-BTC.
    \end{enumerate}
    \item \textbf{Hạn chế kỹ thuật:} (1) Khi đưa một miền tài liệu mới vào framework, quá trình để \texttt{Coding Agent} viết và tinh chỉnh \texttt{Parser} vẫn cần chạy test lặp lại ở đầu kỳ để xử lý hết các edge cases định dạng và thỏa mãn 5 output invariants; (2) Chiến lược giữ nguyên vẹn mỗi mục AST thành một \texttt{chunk} tạo ra phân phối độ dài dạng heavy-tailed với P95 đạt 2.518--3.112 ký tự và cực đại 15.726 ký tự, vượt quá context window 256 token của mô hình Cross-Encoder hiện tại nên bước reranking chỉ quan sát được phần đầu của các mục quá dài; và (3) Cấu hình FTS \texttt{simple} không tự động thực hiện sub-token splitting cho các định danh ghép dạng \texttt{camelCase}/\texttt{PascalCase} trong \texttt{tsvector} index, đồng thời quy tắc \texttt{Sparse Match Gating} có thể gán nhãn \texttt{low} cho các query hợp lệ nhưng không có lexical overlap.
    \item \textbf{Hướng phát triển:} (1) Xây dựng công cụ tự động khép kín vòng lặp để \texttt{Coding Agent} tự lấy mẫu một kho tài liệu mới bất kỳ, sinh code \texttt{Parser}, chạy bộ test 5 invariants và tự sửa lỗi cho đến khi nạp thành công toàn bộ kho dữ liệu vào framework; (2) Thay thế tầng Cross-Encoder 256 token bằng mô hình Long-Context Reranker hoặc cơ chế chấm điểm trượt theo sliding window để bao phủ toàn bộ phần thân của các mục tài liệu dài trên 2.500 ký tự; và (3) Bổ sung bước sub-token splitting cho compound identifiers trước khi đưa vào \texttt{tsvector} index cũng như hỗ trợ trích xuất tự động edge candidates cho Knowledge Graph trong pha ingestion.
\end{itemize}
